\documentclass[11pt]{article}
\usepackage[margin=1in]{geometry}
\usepackage{amsmath,amssymb}
\usepackage{booktabs}
\usepackage{listings}
\usepackage{xcolor}
\usepackage[colorlinks=true,linkcolor=blue,citecolor=blue,urlcolor=blue]{hyperref}

\lstset{
  basicstyle=\ttfamily\footnotesize,
  breaklines=true,
  frame=single,
  columns=fullflexible,
  keepspaces=true,
}

\newcommand{\ket}[1]{\lvert #1 \rangle}
\newcommand{\bra}[1]{\langle #1 \rvert}
\newcommand{\tool}{\textsc{cforge-fuzz}}

\title{Phase-Sensitive Differential Fuzzing for Quantum Simulators:\\
Evaluation and a Soundness Regression in the Qiskit Transpiler}
\author{Cleiton Augusto Correa Bezerra\\
\small \texttt{augusto.cleiton@gmail.com}}
\date{July 2026 --- working draft}

\begin{document}
\maketitle

\begin{abstract}
Quantum software stacks are validated mostly by benchmarks that score a
circuit on its output distribution: Quantum Volume, cross-entropy
benchmarking, mirror circuits. These benchmarks share a blind spot. A
companion result proves that any score computed from computational-basis
measurement statistics is invariant under complex conjugation of every
gate, so an entire class of gate-convention bug passes them by
construction. We present \tool{}, a differential and metamorphic fuzzer
whose oracle compares statevector \emph{amplitudes} modulo a single
global phase rather than output distributions, and which therefore sees
the bugs the benchmarks cannot. We run it against six production stacks:
Qiskit, tket, PennyLane, Cirq, and the quantrs2 Rust simulator. Across
roughly 76{,}000 random circuits the tool reports one previously unknown
soundness bug, a regression introduced in Qiskit~2.5.0 where the
\texttt{CommutativeCancellation} pass deletes a
$\sqrt{X}^\dagger\!\cdot\sqrt{X}^\dagger$ pair that equals $X$, flipping
measurement outcomes at optimization level~2 and above. The same class
of optimization in the other four stacks handles the analogous case
correctly. Our main methodological finding is that an amplitude-level
metamorphic oracle is only sound once it quotients the transformations a
target is licensed to perform: global phase, layout permutation,
approximate synthesis, and finite numerical precision. In each case the
minimally shrunk counterexample identifies which quotient a false
positive came from.
\end{abstract}

\section{Introduction}

A quantum program is a circuit, and the tools that run it --- simulators,
transpilers, optimizers --- are ordinary software with ordinary bugs. The
question is how those bugs are caught. In practice the dominant check is
a benchmark: the developer runs a known circuit, samples measurement
outcomes, and compares the resulting histogram against an expected one.
Quantum Volume~\cite{cross2019}, linear cross-entropy benchmarking
(XEB)~\cite{arute2019}, mirror-circuit benchmarking~\cite{proctor2022},
and the QED-C application suite all follow this pattern.

This kind of check has a structural weakness that, to our knowledge, has
not been stated in the testing literature. Every one of these benchmarks
scores a device on a function of computational-basis measurement
probabilities, that is, on $\lvert\bra{x}C\ket{\psi_0}\rvert^2$ for basis
states $x$ and a real initial state $\psi_0$. A companion paper proves
that this quantity is unchanged if every gate $U$ in the circuit is
replaced by its complex conjugate $U^*$. A simulator that implements a
consistent $i \mapsto -i$ convention error therefore earns an identical
score to a correct one. The bug is real, it changes amplitudes, and no
sampling benchmark can see it.

We take this as a design constraint for a fuzzer. If the goal is to
catch convention bugs, the oracle must look at amplitudes and not at
probabilities. We built \tool{} around an amplitude oracle: it compares
two statevectors up to a single global phase, which is the only genuinely
unobservable degree of freedom. On top of the oracle sits a typed random
circuit generator that oversamples the gates where convention bugs live,
and a delta-debugging shrinker that reduces any divergence to a minimal
witness.

The contribution of this paper is the evaluation. We ran \tool{} against
six widely used stacks and report what it found. The headline is a single
bug, but a real one: a regression in the current release of Qiskit whose
transpiler silently miscompiles a two-gate identity and changes
measurement results. We report it, give its root cause, and trace the
release in which it appeared. We also report the negative results,
because on this class of test they carry information: four mature
optimizers handle the analogous case correctly across tens of thousands
of circuits, which is what lets us attribute the Qiskit failure to a
specific regression rather than to an inherent difficulty.

The second contribution is methodological. Working at amplitude precision
forces a discipline that distribution-level testing can skip. An
optimizer is allowed to change a circuit in several ways that do not
change its meaning, and a naive amplitude comparison flags all of them.
We identify four such licensed transformations, show that ignoring any of
them floods the tool with false positives, and observe that the shrunk
witness itself diagnoses which transformation was missed.

\section{Background and problem}

\paragraph{Gate conventions.}
The OpenQASM~3 standard library fixes a matrix for each named gate. Two
implementations that both claim to follow it can still disagree on a
sign. The rotation $R_z(\lambda) = \mathrm{diag}(e^{-i\lambda/2},
e^{+i\lambda/2})$ is a common point of divergence; some libraries
implement it with the opposite sign. On a basis state this is only a
global phase and is invisible, which is why such bugs survive casual
testing. On a superposition it changes relative phases, and on a
two-qubit entangling circuit it changes amplitudes that a careful oracle
can detect.

\paragraph{Why benchmarks miss them.}
Let $C^*$ denote the circuit with every gate conjugated. For any real
$\psi_0$ and basis state $x$,
\[
\lvert \bra{x} C^* \ket{\psi_0} \rvert^2
= \lvert \overline{\bra{x} C \ket{\psi_0}} \rvert^2
= \lvert \bra{x} C \ket{\psi_0} \rvert^2 ,
\]
so any score that is a function of these probabilities is identical for
$C$ and $C^*$. The companion paper extends this to adaptive protocols
with mid-circuit measurement, and derives it as a corollary for Quantum
Volume, XEB, mirror circuits, and randomized benchmarking. The practical
consequence for testing is the starting point of this work: to catch a
conjugation-class bug, an oracle must compare amplitudes.

\paragraph{Problem statement.}
We want a fuzzer for quantum circuit tooling that (i) uses an oracle
strong enough to detect amplitude-level divergences, (ii) does not report
the many meaning-preserving rewrites that a correct optimizer is entitled
to make, and (iii) turns each true divergence into a witness small enough
to act on. The rest of the paper describes such a tool and what it found.

\section{Approach}

\subsection{The amplitude oracle}
Given two statevectors $a, b \in \mathbb{C}^{2^n}$ for the same circuit,
the oracle computes
\[
N1(a,b) = \min_{\varphi}\ \max_i \bigl\lvert e^{i\varphi} a_i - b_i \bigr\rvert ,
\]
minimized in closed form at $\varphi = -\arg\langle a \vert b\rangle$.
This reports zero exactly when $a$ and $b$ differ only by a global phase,
and grows with any observable amplitude difference. We also compute a
probability distance $N2 = \max_i \lvert\,\lvert a_i\rvert^2 -
\lvert b_i\rvert^2\rvert$ and a single-qubit observable distance $N3$. The
three form a strict hierarchy: a divergence at $N1$ with agreement at
$N2$ is exactly a conjugation-class bug, invisible to every sampling
benchmark. The quantrs2 $R_z$ bug is of this kind, with $N1 = 0.88$ and
$N2 = 0$.

\subsection{Generator}
Circuits are drawn from the OpenQASM standard gate set with a weighted
distribution. The phase family ($R_z$, $P$, controlled phase) and the
discrete complex gates ($S$, $S^\dagger$, $T$, $T^\dagger$, $\sqrt{X}$,
$\sqrt{X}^\dagger$) are weighted at five times the plain Pauli and
structural gates. The reason is specific to convention bugs: such a bug
becomes observable only when a mis-implemented gate meets a correctly
implemented complex gate, so the generator biases toward circuits that
mix the two.

\subsection{Shrinker}
A divergence in a twelve-gate circuit is not actionable. The shrinker
applies delta debugging: it removes chunks of the circuit, keeping any
removal that preserves the divergence, then finishes with a greedy
single-gate pass until the witness is one-minimal, meaning no single
remaining gate can be dropped. The output is a witness of a handful of
gates together with an OpenQASM~2 reproducer. In our campaigns every
Qiskit transpiler divergence shrank to the same three-gate core.

\subsection{Two testing modes}
We use the oracle in two ways. In \emph{differential} mode we run the
same circuit through two implementations of the same semantics, for
example a reference statevector simulator against a third-party one, and
flag any $N1$ divergence. In \emph{metamorphic} mode we run one
implementation before and after a transformation that is supposed to
preserve meaning, such as an optimization pass, and flag any change in
the state. The Qiskit bug was found in metamorphic mode.

\section{Implementation}

The core generator, oracle, and shrinker are implemented in Rust
(\tool{}, roughly 1{,}500 lines) over a shared circuit representation, and
drive a native statevector simulator as the reference. Each external
target is reached through a thin adapter: for Qiskit, tket, PennyLane,
and Cirq we emit the generated circuit into the framework's own API and
read back its statevector or its post-optimization circuit. The adapters
are small Python drivers that share the generator's gate weights and the
oracle's $N1$ definition, so that a divergence means the same thing
across targets. All campaigns are seeded and reproducible, and each
confirmed divergence is stored as a JSON record with the framework
version, the seed, the minimal gate list, and the reproducer.

\section{Evaluation}

We ran one campaign per target. Table~\ref{tab:campaigns} lists the
targets, the mode, the number of random circuits, and the result.

\begin{table}[h]
\centering
\small
\begin{tabular}{llrl}
\toprule
Target & Mode & Circuits & Result \\
\midrule
quantrs2 0.2.0 & differential & 10{,}000 & bug ($R_z$ sign, known) \\
Qiskit 2.5.0 Aer & differential & 3{,}000 & clean \\
Qiskit 2.5.0 transpiler & metamorphic & 8{,}000 & \textbf{bug} (\S\ref{sec:case}) \\
tket 2.18 (4 passes) & metamorphic & 16{,}000 & clean \\
PennyLane 0.42 (lightning) & differential & 3{,}000 & clean \\
PennyLane 0.42 (5 transforms) & metamorphic & 20{,}000 & clean \\
Cirq 1.5 (4 transformers) & metamorphic & 16{,}000 & clean \\
\bottomrule
\end{tabular}
\caption{Campaign results. One previously unknown soundness bug across
roughly 76{,}000 random circuits. The mature frameworks are otherwise
correct on this gate set.}
\label{tab:campaigns}
\end{table}

The yield is one bug. We think this is the honest and expected shape of
the result rather than a disappointing one. Differential testing of
mature software tends to find few but high-value defects; the canonical
example is Csmith, whose value was never the count but the demonstration
that a whole class of compiler bug was reachable~\cite{csmith}. Our
negative results are load-bearing: they show that the amplitude oracle
does not simply flag noise everywhere, and they isolate the Qiskit
failure as specific rather than generic.

\section{Case study: a soundness regression in \texttt{CommutativeCancellation}}
\label{sec:case}

The property under test is that a transpiler preserves the circuit's
unitary up to global phase. The oracle compares $\mathrm{Op}(C)$ with
$\mathrm{Op}(\mathrm{transpile}(C))$, after undoing the transpiler's
declared qubit permutation and disabling approximate synthesis
(Section~\ref{sec:quotients} explains why both are necessary). Every
divergence shrank to the following three-gate core.

\begin{lstlisting}[language=Python]
from qiskit import QuantumCircuit
from qiskit.transpiler import PassManager
from qiskit.transpiler.passes import CommutativeCancellation
from qiskit.quantum_info import Operator

qc = QuantumCircuit(2)
qc.sxdg(1); qc.sx(0); qc.sxdg(1)

out = PassManager([CommutativeCancellation()]).run(qc)
print(out)                                 # only sx(0) remains
print(Operator(qc).equiv(Operator(out)))   # False
\end{lstlisting}

The two $\sqrt{X}^\dagger$ gates on qubit~1 are separated by a gate on
qubit~0 and are cancelled as if they were an inverse pair. They are not.
$\sqrt{X}^\dagger \cdot \sqrt{X}^\dagger = R_x(-\pi)$, which equals $X$ up
to a global phase, not the identity. At optimization level~2 and above,
where the pass runs, the transpiled circuit measures qubit~1 as $0$ with
certainty where the original circuit measures it as $1$.

\paragraph{Root cause.}
The pass sorts gates into cancellable families. In
\texttt{commutative\_cancellation.py} the relevant set is
\[
\texttt{\_x\_rotations} = \{\texttt{"x"}, \texttt{"rx"}, \texttt{"sx"}, \texttt{"sxdg"}\}.
\]
Two identical $\sqrt{X}^\dagger$ gates are treated as mutual inverses.
This is correct for \texttt{x} (self-inverse) but wrong for the
square-root gates, whose square is $X$. Reading the source across
releases, the set was $\{\texttt{"x"},\texttt{"rx"}\}$ from Qiskit~1.4.0
through 2.4.0 and gained \texttt{sx} and \texttt{sxdg} in 2.5.0. The bug
is a regression in the current release. The adjacent case, two
$\sqrt{X}^\dagger$ with nothing between them, is folded correctly into a
single gate; the fault needs an interleaved operation on another qubit,
which is why hand tests miss it and random circuits find it. We reported
the bug upstream as Qiskit issue~16594.

\paragraph{Contrast with the other stacks.}
The same class of pass elsewhere handles the analogous
$\sqrt{X}\cdot\sqrt{X} = X$ correctly. PennyLane's
\texttt{cancel\_inverses}, \texttt{merge\_rotations}, and
\texttt{single\_qubit\_fusion} are clean over 20{,}000 circuits; tket's
\texttt{RemoveRedundancies}, \texttt{CliffordSimp}, and
\texttt{SynthesiseTket} are clean over 16{,}000. The failure is specific
to the Qiskit 2.5.0 regression, not an inherent hazard of the
transformation.

\section{Quotienting licensed transformations}
\label{sec:quotients}

A metamorphic oracle for an optimizer is sound only if it compares
circuits up to exactly the changes the optimizer is allowed to make.
Ignore one of those changes and the tool drowns in false positives.
We met four such transformations. In each case the shrunk witness told
us which one we had missed.

\begin{table}[h]
\centering
\small
\begin{tabular}{lll}
\toprule
Licensed transformation & False-positive signature & Remedy \\
\midrule
Global phase & $N1$ constant, $N2 = 0$ & minimize over $\varphi$ \\
Layout permutation & witness contains a \texttt{swap} & compose with layout \\
Approximate synthesis & lone small-angle gate & disable approximation \\
Numerical precision & uniform $N1 \approx 10^{-7}$ everywhere & force double precision \\
\bottomrule
\end{tabular}
\caption{The four quotients and how the minimal witness reveals each.}
\label{tab:quotients}
\end{table}

Two of these are worth spelling out. Comparing raw statevectors against
the Qiskit transpiler produced 119 false positives in 3{,}000 circuits,
and every one of their minimal witnesses contained a \texttt{swap}. The
transpiler is entitled to permute qubits and records the permutation in
its layout; once the oracle composed with that layout, the false
positives vanished and the real bug remained. In Cirq, the default
simulator returns single-precision states. Every transformer, including
the ones that are exact by construction, reported the same worst-case
$N1 \approx 1.3 \times 10^{-7}$, which is the float32 epsilon. A genuine
soundness bug in one pass does not produce an identical error floor in an
unrelated pass; a uniform floor is a property of the harness, not the
target. Forcing double precision removed all of it.

This discipline is the concrete price of working at amplitude precision,
and also its payoff. A distribution-level oracle never has to name these
quotients because it is too coarse to see them, and it is correspondingly
too coarse to see the conjugation-class bug that motivates the whole
approach. By being explicit about the algebraic and numerical gauge, the
tool can separate a genuine $X$-dropping bug from a $10^{-7}$ rounding
artifact, and the shrinker turns the separation into a diagnosis.

\section{Related work}

The closest prior work is metamorphic and differential testing of quantum
platforms. MorphQ applies metamorphic relations to Qiskit and compares
output distributions~\cite{morphq}; QDiff runs equivalent circuits across
backends and likewise compares measurement statistics~\cite{qdiff}. Both
operate at the level of probabilities, which is exactly the level at
which the conjugation-class bug is invisible. Our oracle differs in
kind, not degree: it compares amplitudes modulo global phase, and it is
motivated by a proof that the distribution level cannot detect the target
class. The cost of that precision is the quotienting discipline of
Section~\ref{sec:quotients}, which distribution-level tools do not need
and cannot exploit.

The broader inspiration is compiler fuzzing, where Csmith established that
random differential testing finds real miscompilations in mature
compilers~\cite{csmith}, and later work formalized the metamorphic style
of oracle~\cite{metamorphic}. On the quantum side, the benchmarks we
argue past --- Quantum Volume~\cite{cross2019}, XEB~\cite{arute2019},
mirror circuits~\cite{proctor2022} --- were designed to characterize
noisy hardware, not to certify gate conventions in software, which is
why their blind spot has gone unremarked rather than being a known flaw.

\section{Threats to validity}

Our campaigns cover up to five qubits and fourteen gates. Deeper circuits
may exercise passes that our generator underweights, so absence of a
divergence is evidence of correctness on the sampled distribution, not a
proof. ``Clean'' means no $N1$ divergence above $10^{-9}$ after the four
quotients; a subtler bug below that threshold, or one requiring a gate we
did not generate, would be missed. The confirmed bug count is one, which
is small; we have argued that this is the expected yield for this style
of testing on mature tools, but a reader who weights bug count heavily
should discount the result accordingly. Finally, the bug report is filed
but not yet triaged upstream at the time of writing, and our root-cause
reading of the source, though supported by the reproducer, is ours rather
than the maintainers'.

\section{Conclusion}

Sampling benchmarks cannot certify gate conventions, and a companion
result makes that precise. Building a fuzzer around an amplitude oracle
closes the gap, at the cost of a quotienting discipline that we have made
explicit. Run against six production stacks, the tool found a real
soundness regression in the current Qiskit transpiler and confirmed the
correctness of four other optimizers on the same class of input. The bug
is one; the method, and the proof that explains why the class was
overlooked, are the durable part.

\begin{thebibliography}{9}
\bibitem{cross2019} A.~Cross, L.~Bishop, S.~Sheldon, P.~Nation,
J.~Gambetta, \emph{Validating quantum computers using randomized model
circuits}, Phys.\ Rev.\ A \textbf{100}, 032328 (2019).
\bibitem{arute2019} F.~Arute \emph{et al.}, \emph{Quantum supremacy using
a programmable superconducting processor}, Nature \textbf{574}, 505
(2019).
\bibitem{proctor2022} T.~Proctor \emph{et al.}, \emph{Measuring the
capabilities of quantum computers}, Nat.\ Phys.\ \textbf{18}, 75 (2022).
\bibitem{morphq} M.~Paltenghi and M.~Pradel, \emph{MorphQ: Metamorphic
testing of the Qiskit quantum computing platform}, ICSE 2023.
\bibitem{qdiff} J.~Wang \emph{et al.}, \emph{QDiff: Differential testing
of quantum software stacks}, ASE 2021.
\bibitem{csmith} X.~Yang, Y.~Chen, E.~Eide, J.~Regehr, \emph{Finding and
understanding bugs in C compilers}, PLDI 2011.
\bibitem{metamorphic} T.~Y.~Chen \emph{et al.}, \emph{Metamorphic
testing: A review of challenges and opportunities}, ACM Comput.\ Surv.\
\textbf{51}(1), 2018.
\end{thebibliography}

\end{document}
